\documentclass[12pt]{article}
\usepackage[utf8]{inputenc}

\usepackage{mystyle}

\title{Complex Curves}
\author{Joseph Sullivan}
\date{July 2025}

\DeclareMathOperator{\Kom}{Kom}
\DeclareMathOperator{\Cyl}{Cyl}
\DeclareMathOperator{\Stab}{Stab}
\DeclareMathOperator{\Slice}{Slice}
\DeclareMathOperator{\chern}{ch}
\DeclareMathOperator{\todd}{td}
\DeclareMathOperator{\Spin}{Spin}

\DeclareMathOperator{\ext}{ext}
%\DeclareMathOperator{\hom}{hom}

\newcommand{\cQ}{\mathcal{Q}}
\newcommand{\cZ}{\mathcal{Z}}
\newcommand{\LCom}{\mathcal{LC}}

\newcommand{\rddots}{\text{\reflectbox{$\ddots$}}}


\begin{document}

\maketitle
\tableofcontents


\section{When is $\omega_C$ Very Ample?}

\begin{defn}
    A \textbf{curve} is a 1 dimensional complex projective manifold (smooth variety).
\end{defn}

\begin{defn}
    The \textbf{genus} of a curve $C$ is
    \[g = h^0(C,\omega_C) = h^1(C, \cO_C).\]
\end{defn}

\begin{defn}
    A curve $C$ is \textbf{hyperelliptic} if it admits a degree 2 map $C\to \P^1$.
\end{defn}

\begin{thm}[Riemann Roch for Curves]
    Let $L$ be a line bundle on $C$.
    \[\chi(L) = \deg L - g + 1.\]
\end{thm}

Our first main theorem is about when the canonical bundle $\omega_C$ gives a map to projective space. Clearly we need $g\geq 1$ because this is the degree of $\omega_C$, and this is enough to be ample by Riemann Roch/closed embedding test for curves. However, we want to know when $\omega_C$ is \textit{very ample}. We will work towards the following result.

\begin{thm}
    If $g\geq 1$, either $\omega_C$ is very ample OR $C$ is hyperelliptic (cannot have both).
\end{thm}

To prove this, we will need bits of Riemann Hurwitz and clever applications of Riemann Roch.

\subsection{Riemann Hurwitz}

\begin{defn}
    Let $f:X\to Y$ be a dominant map of curves. (A lot of these definitions work more generally).
    \begin{itemize}
        \item The \textbf{degree} of $f$ is the degree of the extension of function fields $k(X)/k(Y)$. This will turn out to be the general number of preimages of a closed point.
        \item For $P\in X$, the \textbf{ramification index} $e_P$ is the valuation of the pullback of a local parameter $t\in \cO_{f(P)}$. Roughly, $t$ parametrizes $Y$, and we're checking how much $f$ vanishes in the $t$ direction to see how many sheets come together.
        \item If $e_P > 1$, say $f$ is \textbf{ramified} at $P$ and $f(P)$ is a branch point.
        \item Define $f^*: \Div Y \to \Div X$ by
        \[f^*(Q) = \sum_{P\mapsto Q} e_P \cdot P\]
        and extending by linearity. This is compatible with $f^*$ on $\Pic$.
    \end{itemize}
\end{defn}

\begin{prop}
    There is a short exact sequence (conormal sequence, plus injectivity)
    \[0\to f^* \Omega_Y \to \Omega_X \to \Omega_{X/Y} \to 0.\]
\end{prop}

\begin{prop}
    $\Omega_{X/Y}$ is a torsion sheaf supported on the ramification points of $f$, and
    \[\mathrm{length}(\Omega_{X/Y})_P = e_P -1,\]
    hence if we define the \textbf{ramification divisor} by
    \[R = \sum_{P \text{ ramification point}} (e_P-1)\cdot P,\]
    we get $\Omega_{X/Y} = \cO_R$.
\end{prop}

\begin{thm}[Riemann Hurwitz]
    $K_X \sim f^* K_Y + R$, hence
    \[2g(X) - 2 = (\deg f)(2g(Y)-2) + \sum_{P\in X}(e_P-1).\]
\end{thm}

\begin{cor}
    Let $\pi: C\to \P^1$ be a double cover. Then, $\pi$ has $2g+2$ branch points.
\end{cor}
\begin{proof}
    By Riemann-Hurwitz, we get
    \[2g-2 = 2(0-2)+ \sum_{P\in X} e_P-1,\]
    and $e_P \leq 2$ {\color{red} because double cover(?)}. So, there are exactly $2g+2$ ramification points.
\end{proof}


% \nocite{*} % Include all references in refs.bib regardless of whether they are referenced or not
% \bibliographystyle{plain} % We choose the "plain" reference style
% \bibliography{2025summer/quadric} % Entries are in the refs.bib file

\end{document}
